\begin{proof}[Proof of \cref{obj:lem:full-marked-likelihood}]
\begin{enumerate}
\item Define the perturbation amplitude and Gaussian density by
\[
\eta_{\mathrm{pert}}
=\kappa\left(\frac{b}{m^2}\right)^\beta,
\qquad
\varphi_\sigma(t)
=\frac{\exp(-t^2/(2\sigma^2))}{\sqrt{2\pi}\,\sigma},
\quad t\in\mathbb R.
\]
Since \(m\ge2\), \(0<b\le1\), and \(\beta>0\),
\[
0<\frac{b}{m^2}\le1,
\qquad
0<\left(\frac{b}{m^2}\right)^\beta\le1,
\qquad
0<\eta_{\mathrm{pert}}\le\frac1{100}\le\frac14.
\]
The polynomial bound in \cref{obj:lem:legal-cancellation}, obtained from the identities of \cref{obj:lem:classical-legendre-facts}, and the extension in \cref{obj:def:legal-extension} give
\[
|v_{b,m}(x)|\le1,
\qquad
\frac14\le\frac12\pm\eta_{\mathrm{pert}}v_{b,m}(x)\le\frac34,
\quad x\in[-1,1].
\]
Consequently, the conditional mass of either value of \(Y^1\) under the positive alternative is at most three times its conditional mass under the negative alternative. The score, \(Y^0\), and error distributions in \cref{obj:def:alternatives} are common to the two laws. Integrating the conditional mass comparison and then applying the observation map therefore gives
\[
P^+_{b,m}\le3P^-_{b,m},
\qquad
(P^+_{b,m})_{\mathrm{obs}}
\le3(P^-_{b,m})_{\mathrm{obs}}.
\]
In particular,
\[
(P^+_{b,m})_{\mathrm{obs}}
\ll(P^-_{b,m})_{\mathrm{obs}}.
\]
Define the likelihood ratio, up to almost-everywhere equality, by
\[
\mathsf{LR}(w,d,y)
=
\frac{d(P^+_{b,m})_{\mathrm{obs}}}
     {d(P^-_{b,m})_{\mathrm{obs}}}(w,d,y),
\qquad
(w,d,y)\in\mathbb R\times\{0,1\}^2.
\]
Domination implies
\[
0\le\mathsf{LR}\le3,
\qquad
(\mathsf{LR}-1)^2\le4
\quad
(P^-_{b,m})_{\mathrm{obs}}\text{-almost everywhere}.
\]
Both observed laws are probability laws, so the squared likelihood-ratio deviation is integrable. The extended divergence consequently has the finite integral representation
\[
\chi^2\!\left((P^+_{b,m})_{\mathrm{obs}},
             (P^-_{b,m})_{\mathrm{obs}}\right)
=
\int(\mathsf{LR}-1)^2\,d(P^-_{b,m})_{\mathrm{obs}}.
\]
% lean: CausalSmith.Stat.RdTruesideNoiseFrontier.altLaw_observed_absolutelyContinuous
% lean: CausalSmith.Stat.RdTruesideNoiseFrontier.altLaw_observed_likelihood_square_integrable
% lean: CausalSmith.Stat.RdTruesideNoiseFrontier.chiSqExt_eq_ofReal_chiSqDiv

\item The positivity of \(\varphi_\sigma\), the bound \(b\le1\), and \cref{obj:lem:legal-cancellation} imply, for every \(w\in\mathbb R\),
\[
q_\sigma(w)>0,
\qquad
\begin{aligned}
|u_{b,m,\sigma}(w)|
&\le\int_0^b |g_m(x/b)|\varphi_\sigma(w-x)\,dx\\
&\le\int_0^b\varphi_\sigma(w-x)\,dx
\le q_\sigma(w).
\end{aligned}
\]
Gaussian normalization and integration over the unit latent interval give
\[
\int_{\mathbb R}q_\sigma(w)\,dw
=
\int_0^1\left(\int_{\mathbb R}\varphi_\sigma(w-x)\,dw\right)dx
=1.
\]
Thus
\[
0\le
\frac{u_{b,m,\sigma}(w)^2}{q_\sigma(w)}
\le q_\sigma(w),
\]
which proves the asserted Lebesgue integrability.

To compute the full observed likelihood, let
\[
\mathsf p_\pm(w,d,y)
=
\frac{d(P^\pm_{b,m})_{\mathrm{obs}}}
     {d(\text{Lebesgue measure}\otimes\text{counting measure on }\{0,1\}^2)}
     (w,d,y).
\]
Versions defined for every \(w\in\mathbb R\) and \(d,y\in\{0,1\}\) are
\[
\begin{aligned}
\mathsf p_\pm(w,1,y)
&=\frac12\int_0^1
 \left(\frac12\pm(2y-1)\eta_{\mathrm{pert}}v_{b,m}(x)\right)
 \varphi_\sigma(w-x)\,dx,\\
\mathsf p_\pm(w,0,y)
&=\frac14\int_{-1}^0\varphi_\sigma(w-x)\,dx.
\end{aligned}
\]
These formulas follow by conditioning on the uniform score and using the independent Gaussian error and the conditional Bernoulli probabilities in \cref{obj:def:alternatives}. Splitting the first integral at \(b\) and using \cref{obj:def:legal-extension} gives
\[
\begin{aligned}
\int_0^1v_{b,m}(x)\varphi_\sigma(w-x)\,dx
&=\int_0^bg_m(x/b)\varphi_\sigma(w-x)\,dx
  +\int_b^1 0\,dx\\
&=u_{b,m,\sigma}(w).
\end{aligned}
\]
This identity also covers \(b=1\), when the second integral is zero. Hence
\[
\begin{aligned}
\mathsf p_\pm(w,1,1)
&=\frac{q_\sigma(w)}4
  \pm\frac{\eta_{\mathrm{pert}}u_{b,m,\sigma}(w)}2,\\
\mathsf p_\pm(w,1,0)
&=\frac{q_\sigma(w)}4
  \mp\frac{\eta_{\mathrm{pert}}u_{b,m,\sigma}(w)}2.
\end{aligned}
\]
The untreated densities agree, and all four cell densities are positive. Therefore the likelihood ratio has the version
\[
\mathsf{LR}(w,d,y)
=
\begin{cases}
\mathsf p_+(w,1,y)/\mathsf p_-(w,1,y),&d=1,\\
1,&d=0.
\end{cases}
\]
Integrating with the negative-law density and summing over the binary cells yields
\[
\begin{aligned}
&\chi^2\!\left((P^+_{b,m})_{\mathrm{obs}},
              (P^-_{b,m})_{\mathrm{obs}}\right)\\
&\quad=
\int_{\mathbb R}\sum_{y\in\{0,1\}}
\frac{\bigl(\mathsf p_+(w,1,y)-\mathsf p_-(w,1,y)\bigr)^2}
     {\mathsf p_-(w,1,y)}\,dw.
\end{aligned}
\]
The untreated contributions vanish because their likelihood ratio is one.
% lean: CausalSmith.Stat.RdTruesideNoiseFrontier.markedConvolution_square_div_integrable
% lean: CausalSmith.Stat.RdTruesideNoiseFrontier.altLaw_observed_rnDeriv
% lean: CausalSmith.Stat.RdTruesideNoiseFrontier.altLikelihoodRatio_sq_integral_eq

\item For every \(w\in\mathbb R\),
\[
|\eta_{\mathrm{pert}}u_{b,m,\sigma}(w)|
\le\frac{q_\sigma(w)}4,
\qquad
\frac{q_\sigma(w)}4
\pm\frac{\eta_{\mathrm{pert}}u_{b,m,\sigma}(w)}2
\ge\frac{q_\sigma(w)}8.
\]
The two treated-cell differences are
\[
\begin{aligned}
\mathsf p_+(w,1,1)-\mathsf p_-(w,1,1)
&=\eta_{\mathrm{pert}}u_{b,m,\sigma}(w),\\
\mathsf p_+(w,1,0)-\mathsf p_-(w,1,0)
&=-\eta_{\mathrm{pert}}u_{b,m,\sigma}(w).
\end{aligned}
\]
Thus their squared likelihood contributions satisfy
\[
\begin{aligned}
&\sum_{y\in\{0,1\}}
\frac{\bigl(\mathsf p_+(w,1,y)-\mathsf p_-(w,1,y)\bigr)^2}
     {\mathsf p_-(w,1,y)}\\
&\quad=
\frac{(\eta_{\mathrm{pert}}u_{b,m,\sigma}(w))^2}
     {q_\sigma(w)/4-\eta_{\mathrm{pert}}u_{b,m,\sigma}(w)/2}
+
\frac{(\eta_{\mathrm{pert}}u_{b,m,\sigma}(w))^2}
     {q_\sigma(w)/4+\eta_{\mathrm{pert}}u_{b,m,\sigma}(w)/2}\\
&\quad\le
2\,\frac{(\eta_{\mathrm{pert}}u_{b,m,\sigma}(w))^2}{q_\sigma(w)/8}
=
16\eta_{\mathrm{pert}}^2
\frac{u_{b,m,\sigma}(w)^2}{q_\sigma(w)}.
\end{aligned}
\]
The right-hand side is integrable by the preceding step, and the nonnegative left-hand side is dominated by it. Integrating and using
\[
\eta_{\mathrm{pert}}^2
=\kappa^2\left(\frac{b}{m^2}\right)^{2\beta}
\]
proves the first inequality.
% lean: CausalSmith.Stat.RdTruesideNoiseFrontier.witnessTreatedCellConvolution_integral_bound

\item Define the centered moments for every integer \(j\ge0\) by
\[
\zeta_j
=
\int_0^b g_m(x/b)(x-b/2)^j\,dx.
\]
The change of variable \(x=bt\) gives
\[
\zeta_j
=b\int_0^1g_m(t)(bt-b/2)^j\,dt.
\]
For \(0\le j\le m-2\), the second factor is a polynomial of degree at most \(j\). Expanding it into monomials and applying the cancellation identities of \cref{obj:lem:legal-cancellation} therefore gives
\[
\zeta_j=0,\qquad 0\le j\le m-2.
\]
The same cited result and scaling give
\[
\int_0^b g_m(x/b)^2\,dx
=b\int_0^1g_m(t)^2\,dt
\le\frac b{N_m}.
\]
Also, \(|x-b/2|\le b/2\) on \([0,b]\), so
\[
\int_0^b\bigl((x-b/2)^j\bigr)^2\,dx
\le b(b/2)^{2j}.
\]
Since \(N_m=3m^2+4m+1>0\), Cauchy--Schwarz yields
\[
\begin{aligned}
\zeta_j^2
&\le
\left(\int_0^b g_m(x/b)^2\,dx\right)
\left(\int_0^b\bigl((x-b/2)^j\bigr)^2\,dx\right)\\
&\le\frac{b^2}{N_m}(b/2)^{2j}.
\end{aligned}
\]
Dividing by the positive variance and factorial factors gives
\[
0\le\frac{\zeta_j^2}{\sigma^{2j}j!}
\le
\frac{b^2}{N_m}
\begin{cases}
0,&0\le j<m-1,\\
\lambda_{b,\sigma}^{\,j}/j!,&j\ge m-1,
\end{cases}
\]
because
\[
\frac{(b/2)^{2j}}{\sigma^{2j}}
=\lambda_{b,\sigma}^{\,j}.
\]
The nonnegative exponential series satisfies
\[
\sum_{j=0}^{\infty}\frac{\lambda_{b,\sigma}^{\,j}}{j!}
=e^{\lambda_{b,\sigma}}<\infty.
\]
It follows by comparison that both the retained factorial series and the centered-moment series converge, and
\[
\sum_{j=0}^{\infty}\frac{\zeta_j^2}{\sigma^{2j}j!}
\le
\frac{b^2}{N_m}
\sum_{j=m-1}^{\infty}\frac{\lambda_{b,\sigma}^{\,j}}{j!}.
\]
% lean: CausalSmith.Stat.RdTruesideNoiseFrontier.centered_block_moment_series_bound
% lean: CausalSmith.Stat.RdTruesideNoiseFrontier.likelihoodTail_summable

\item For \(w\in\mathbb R\) and \(x\in[0,b]\), define
\[
\xi_w(x)
=
\frac{-(w-x)^2+(w-b/2)^2+b^2/12}{2\sigma^2}.
\]
Expansion of the numerator and elementary polynomial integration give
\[
\begin{aligned}
\int_0^b\xi_w(x)\,dx
&=\frac1{2\sigma^2}
  \int_0^b\left(2wx-x^2-wb+\frac{b^2}{3}\right)dx\\
&=\frac{wb^2-b^3/3-wb^2+b^3/3}{2\sigma^2}
=0.
\end{aligned}
\]
Since \(e^{\xi_w(x)}\ge1+\xi_w(x)\),
\[
\int_0^b e^{\xi_w(x)}\,dx\ge b.
\]
The Gaussian formula factors as
\[
\varphi_\sigma(w-x)
=
e^{-b^2/(24\sigma^2)}
\varphi_\sigma(w-b/2)e^{\xi_w(x)}.
\]
Integrating this identity and restricting the unit-interval convolution to \([0,b]\) consequently gives, for every real \(w\),
\[
q_\sigma(w)
\ge\int_0^b\varphi_\sigma(w-x)\,dx
\ge
b e^{-b^2/(24\sigma^2)}\varphi_\sigma(w-b/2).
\]
All denominators are positive, so
\[
\frac{u_{b,m,\sigma}(w)^2}{q_\sigma(w)}
\le
\frac{e^{b^2/(24\sigma^2)}}b
\frac{u_{b,m,\sigma}(w)^2}{\varphi_\sigma(w-b/2)}.
\]
% lean: CausalSmith.Stat.RdTruesideNoiseFrontier.marked_square_midpoint_comparison

\item Completing the square gives, for \(x,t,w\in\mathbb R\),
\[
\frac{\varphi_\sigma(w-x)\varphi_\sigma(w-t)}
     {\varphi_\sigma(w-b/2)}
=
\exp\!\left(\frac{(x-b/2)(t-b/2)}{\sigma^2}\right)
\varphi_\sigma\bigl(w-(x+t-b/2)\bigr).
\]
Gaussian normalization therefore yields
\[
\int_{\mathbb R}
\frac{\varphi_\sigma(w-x)\varphi_\sigma(w-t)}
     {\varphi_\sigma(w-b/2)}\,dw
=
\exp\!\left(\frac{(x-b/2)(t-b/2)}{\sigma^2}\right).
\]
On the support square \(x,t\in[0,b]\), \cref{obj:lem:legal-cancellation} gives
\[
|g_m(x/b)g_m(t/b)|\le1,
\qquad
\left|\frac{(x-b/2)(t-b/2)}{\sigma^2}\right|
\le\lambda_{b,\sigma}.
\]
Hence the absolute integral of the signed three-variable kernel is bounded by
\[
\begin{aligned}
&\int_0^b\int_0^b\int_{\mathbb R}
\left|
g_m(x/b)g_m(t/b)
\frac{\varphi_\sigma(w-x)\varphi_\sigma(w-t)}
     {\varphi_\sigma(w-b/2)}
\right|\,dw\,dt\,dx\\
&\quad=
\int_0^b\int_0^b
|g_m(x/b)g_m(t/b)|
\exp\!\left(\frac{(x-b/2)(t-b/2)}{\sigma^2}\right)dt\,dx\\
&\quad\le b^2e^{\lambda_{b,\sigma}}<\infty.
\end{aligned}
\]
Multiplying the two integrals defining \(u_{b,m,\sigma}(w)\) gives
\[
\frac{u_{b,m,\sigma}(w)^2}{\varphi_\sigma(w-b/2)}
=
\int_0^b\int_0^b
g_m(x/b)g_m(t/b)
\frac{\varphi_\sigma(w-x)\varphi_\sigma(w-t)}
     {\varphi_\sigma(w-b/2)}\,dt\,dx.
\]
The absolute-integrability bound justifies Fubini and proves integrability of the left-hand side over \(\mathbb R\). It follows that
\[
\begin{aligned}
&\int_{\mathbb R}
\frac{u_{b,m,\sigma}(w)^2}{\varphi_\sigma(w-b/2)}\,dw\\
&\quad=
\int_0^b\int_0^b
g_m(x/b)g_m(t/b)
\exp\!\left(\frac{(x-b/2)(t-b/2)}{\sigma^2}\right)dt\,dx.
\end{aligned}
\]
In particular, the preceding pointwise midpoint comparison can be integrated:
\[
\int_{\mathbb R}\frac{u_{b,m,\sigma}(w)^2}{q_\sigma(w)}\,dw
\le
\frac{e^{b^2/(24\sigma^2)}}b
\int_{\mathbb R}
\frac{u_{b,m,\sigma}(w)^2}{\varphi_\sigma(w-b/2)}\,dw.
\]
% lean: CausalSmith.Stat.RdTruesideNoiseFrontier.marked_square_reference_integral
% lean: CausalSmith.Stat.RdTruesideNoiseFrontier.marked_square_integral_midpoint_comparison

\item Expand the exponential in the compact double integral. For every \(j\ge0\) and \(x,t\in[0,b]\),
\[
\left|
g_m(x/b)g_m(t/b)
\frac{\bigl((x-b/2)(t-b/2)/\sigma^2\bigr)^j}{j!}
\right|
\le\frac{\lambda_{b,\sigma}^{\,j}}{j!}.
\]
The integrals of these absolute values have the summable majorant
\[
\sum_{j=0}^{\infty}
b^2\frac{\lambda_{b,\sigma}^{\,j}}{j!}
=b^2e^{\lambda_{b,\sigma}}.
\]
Thus termwise integration is justified. Each term factors into the product of two centered moments:
\[
\begin{aligned}
&\int_0^b\int_0^b
g_m(x/b)g_m(t/b)
\frac{\bigl((x-b/2)(t-b/2)/\sigma^2\bigr)^j}{j!}\,dt\,dx\\
&\quad=
\frac1{\sigma^{2j}j!}
\left(\int_0^b g_m(x/b)(x-b/2)^j\,dx\right)^2
=\frac{\zeta_j^2}{\sigma^{2j}j!}.
\end{aligned}
\]
Consequently,
\[
\int_{\mathbb R}
\frac{u_{b,m,\sigma}(w)^2}{\varphi_\sigma(w-b/2)}\,dw
=
\sum_{j=0}^{\infty}\frac{\zeta_j^2}{\sigma^{2j}j!}.
\]
Substituting this identity and the centered-moment bound into the integrated midpoint comparison gives
\[
\begin{aligned}
\int_{\mathbb R}\frac{u_{b,m,\sigma}(w)^2}{q_\sigma(w)}\,dw
&\le
\frac{e^{b^2/(24\sigma^2)}}b
\sum_{j=0}^{\infty}\frac{\zeta_j^2}{\sigma^{2j}j!}\\
&\le
\frac{e^{b^2/(24\sigma^2)}}b\,
\frac{b^2}{N_m}
\sum_{j=m-1}^{\infty}\frac{\lambda_{b,\sigma}^{\,j}}{j!}\\
&=
\frac b{N_m}e^{b^2/(24\sigma^2)}
\sum_{j=m-1}^{\infty}\frac{\lambda_{b,\sigma}^{\,j}}{j!}.
\end{aligned}
\]
Multiplication by the nonnegative factor
\(16\kappa^2(b/m^2)^{2\beta}\) proves the second displayed inequality. Absolute continuity, both integrability assertions, and convergence of the factorial series were established above. The likelihood calculation sums both binary outcome cells and integrates over every real proxy value, so it compares the full observed triples.
% lean: CausalSmith.Stat.RdTruesideNoiseFrontier.marked_square_reference_moment_series
% lean: CausalSmith.Stat.RdTruesideNoiseFrontier.marked_square_factorial_tail_bound
% lean: CausalSmith.Stat.RdTruesideNoiseFrontier.full_marked_likelihood
\end{enumerate}
\end{proof}
